\section{Unbounded sequential endpoints: complete proofs}
\label{sec:annex-p15-robust-sequential-endpoints}

\begin{annexitem}{annex:p15.directional.log.integrability}{thm:p15-directional-log-integrability}{The integrable negative part of a log likelihood}
\begin{proof}
Take the finite dominating measure \(\nu=Q+P\), and finite
nonnegative densities \(q,p\). Under \(Q\), \(q>0\) almost surely.
Define \(Z=\log(q/p)\) there, with value \(+\infty\) when \(p=0\).
Thus \(Z=-\infty\) has \(Q\) probability zero. On \(\{0<q<p\}\),
put \(r=q/p\in(0,1)\). Elementary differentiation shows
\(\sup_{0<r<1}r\log(1/r)=1/e\), whence
\[
 EZ^-=\int_{\{0<q<p\}}p\,r\log(1/r)\,d\nu
                         \leq\frac1e\int p\,d\nu=\frac1e.
\]
The expectation of \(Z\) is therefore an unambiguous extended
integral, never a subtraction of two infinities.
Also
\[
 Ee^{-Z}=\int_{\{q>0\}}p\,d\nu=P(q>0)\leq1.
\]

If \(Q(p=0)>0\), then \(Z=\infty\) with positive probability and
\(EZ=\infty\). Otherwise \(Z\) is finite almost surely.
For real \(z\), \(z-1+e^{-z}\geq0\), with equality exactly at zero,
by differentiation. Since \(Ee^{-Z}\leq1\) and \(EZ^-<\infty\),
integration gives
\[
 EZ=1-Ee^{-Z}+E(Z-1+e^{-Z})\geq0.
\]
If equality held, both nonnegative terms would vanish. We would have
\(Z=0\) \(Q\)-almost surely and \(P(q>0)=1\), implying \(P=Q\).
Thus the expectation is strictly positive for distinct laws.
This is precisely \(D(Q\Vert P)\), including its infinite and singular
cases. The description is unchanged by a common dominating measure.

The earlier score order and its quantile construction couple the
oriented score \(Y_i\) under any \(R\in\mathcal H_i\) above \(Z_i\).
Their negative parts consequently satisfy the opposite pointwise
inequality, proving integrability and excluding negative infinity.
Apply the same order to \(\min(Y_i,M)\) and then increase \(M\):
the negative parts are integrable and the positive parts increase,
so the resulting expectation inequality is \(E_RY_i\geq D_i\).
Finally \(e^{-y}\) is nonnegative and decreasing; truncate it first,
apply the order, and then use monotone convergence to obtain
\(E_Re^{-Y_i}\leq E e^{-Z_i}\leq1\).
\end{proof}
\end{annexitem}

\begin{annexitem}{annex:p15.one.sided.truncation}{thm:p15-one-sided-truncation}{Positive truncation, finite stopping and both entropy scales}
\begin{proof}
For \(M>0\), \(Z^{(M)}=\min(Z,M)\) has the same integrable negative
part as \(Z\), while its positive part is bounded. Hence it is
integrable and \(d_M=EZ^{(M)}\uparrow EZ=D\).
Choose \(M\) large enough that \(d_M>0\).
Let \(\tau_t^{(M)}\) be the one-sided hitting time for the independently
truncated increments. For a fixed integer \(n\), finite summation and
independence of the next increment from its preceding history give
\[
 E S^{(M)}_{\tau_t^{(M)}\wedge n}
 =\sum_{j=1}^n E\!\left[Z_j^{(M)}
                 {\bf1}_{\{\tau_t^{(M)}\geq j\}}\right]
 =d_M E(\tau_t^{(M)}\wedge n).
\]
Each finite stopped sum is integrable. Pathwise it is at most \(t+M\):
before a crossing it is less than \(t\), and at a crossing its
last increment is at most \(M\).
Therefore
\(d_M E(\tau_t^{(M)}\wedge n)\leq t+M\).
Monotone convergence in \(n\) gives
\[
 E\tau_t^{(M)}\leq(t+M)/d_M<\infty,
\]
which also proves almost-sure crossing.
Using the same increments before and after truncation,
\(S_n\geq S_n^{(M)}\) for every \(n\). Thus
\(\tau_t\leq\tau_t^{(M)}\) and the stated upper bound follows.
Since \(t>0\), \(\tau_t\geq1\).

If \(D<\infty\), both parts of \(Z\) are integrable.
The already proved \(E\tau_t<\infty\) justifies signed summation:
\[
 \sum_{j\geq1}E\!\left[|Z_j|{\bf1}_{\{\tau_t\geq j\}}\right]
 =E|Z|\,E\tau_t<\infty.
\]
Independence again gives \(ES_{\tau_t}=D E\tau_t\).
Since \(S_{\tau_t}\geq t\), this yields \(m(t)\geq t/D\).
For each fixed permitted \(M\), the upper bound gives
\(\limsup_{t\to\infty}m(t)/t\leq1/d_M\).
Increase \(M\). If \(D<\infty\), the lower and upper bounds squeeze
the limit to \(1/D\); if \(D=\infty\), \(d_M\to\infty\) and the
nonnegative normalized means tend to zero.

If \(s=\Pr(Z=\infty)>0\), the index \(T_\infty\) of the first infinite
increment is geometric with
\(\Pr(T_\infty>n)=(1-s)^n\), so \(ET_\infty=1/s\).
Every \(\tau_t\leq T_\infty\).
For almost every path, \(T_\infty\) is finite and all its earlier
partial sums are finite. As \(t\) increases beyond their maximum,
\(\tau_t=T_\infty\). Thus \(\tau_t\uparrow T_\infty\) and
monotone convergence gives \(m(t)\to1/s\).
If \(s=0\), every fixed finite prefix has a finite maximum almost surely.
Consequently \(\tau_t\uparrow\infty\) as \(t\to\infty\), and
monotone convergence gives \(m(t)\to\infty\).
No assertion about the expectation of an unbounded overshoot was
assumed at any step.
\end{proof}
\end{annexitem}

\begin{annexitem}{annex:p15.hitting.restart}{thm:p15-hitting-restart}{Independence after stopping with a random target}
\begin{proof}
For each path, \(\tau_t\) is nondecreasing in \(t\), hence so is \(m(t)\).
A finite monotone function on \([0,\infty)\) is Borel measurable.
Joint measurability of a hitting time in its boundary and its real-score
path follows from
\[
 \{\tau_t\leq n\}=\{\max_{0\leq j\leq n}S_j\geq t\}
 \qquad(t>0),
\]
an event determined by finitely many measurable sums and comparisons.

Here is the exact restarting argument. On the event \(\{T=n\}\)
for a stopping time \(T\), all the event's information lies in the
first \(n\) increments. Any fixed finite block
\((Z_{n+1},\ldots,Z_{n+k})\) is independent of that information and
has the original product law. Multiplying finite-prefix indicators
and future-block indicators, taking expectations and summing over
\(n\) proves the corresponding factorization at \(T\).
Finite linear combinations and monotone limits extend this from
indicators to all nonnegative measurable functions of a prefix
quantity and a future block. For the hitting-time expectation, use
\(\tau_r=\sum_{k\geq0}{\bf1}_{\{\tau_r>k\}}\): each summand is
jointly measurable in \(r\) and the first \(k\) future increments.
Its finite partial sums increase to the hitting time, so monotone
convergence supplies exactly the infinite-future functional used here.
Equivalently, apply product integration separately on every
\(\{T=n\}\) and sum by Tonelli.
This proves the restart identity without a conditional distribution
on the original observation space.

After \(\tau_t\), the accumulated score is at least \(t\).
Starting a fresh walk which reaches \(v\) guarantees that the old
walk reaches \(t+v\) no later than \(\tau_t+\tau_v\) for that future.
If the old score is already infinite there is no additional wait.
Take expectations using the proved restart property to get
\(m(t+v)\leq m(t)+m(v)\). The definitions at zero cover those
endpoints as well.

For the two-boundary stop \(N\), pathwise \(N\leq\tau_t\), so
\(EN<\infty\) and \(N\) is almost surely finite.
At a correct exit, \(\tau_t=N\).
At a wrong exit, \(S_N\leq-u\) is finite: a positive infinite
increment would have caused a correct exit, and negative infinite
increments have been excluded.
There has been no earlier crossing of \(t\), so the additional
time is exactly the hitting time of \(t-S_N\) by the fresh future.
The boundary map \(t-S_N\) and the jointly measurable hitting-time
map proved above justify product integration with this random boundary.
Taking expectations gives
\[
 m(t)=EN+E[{\bf1}_Wm(t-S_N)].
\]
All terms are nonnegative and the left side is finite, so no
unproved integrability condition is required for this substitution.
\end{proof}
\end{annexitem}

\begin{annexitem}{annex:p15.two.boundary.passage.comparison}{thm:p15-two-boundary-passage-comparison}{Controlling the wrong-exit overshoot by its exponential}
\begin{proof}
Since \(Ee^{-Z}\leq1\), the nonnegative process \(e^{-S_n}\)
has conditional expectation at the next step at most its present value.
For the bounded stop \(N\wedge n\), condition one step at a time:
on \(\{N\leq n\}\) the stopped process stays fixed; otherwise its
next conditional expectation decreases. Induction gives
\(Ee^{-S_{N\wedge n}}\leq1\).
The earlier finite-passage bound proves \(N<\infty\) almost surely.
Fatou's lemma therefore yields \(Ee^{-S_N}\leq1\), also when
a correct exit has score \(+\infty\) and exponential zero.
On \(W\), \(e^{-S_N}\geq e^u\), giving \(\Pr(W)\leq e^{-u}\).

Write \(v=-S_N\geq u\) on \(W\). The restart identity, subadditivity
and the fixed truncation bound give
\[
 m(t)-EN
 =E[{\bf1}_Wm(t+v)]
 \leq \Pr(W)m(t)+\frac1{d_M}E[{\bf1}_W(v+M)].
\]
For \(u\geq1\) and \(v\geq u\), the function
\((v+M)e^{-v}\) is nonincreasing, since its derivative is
\((1-v-M)e^{-v}\leq0\). Thus
\[
 E[{\bf1}_W(v+M)]
 \leq (u+M)e^{-u}E[{\bf1}_W e^v]
 \leq(u+M)e^{-u}.
\]
This proves integrability of the relevant wrong-exit overshoot rather
than imposing it as a hypothesis.
Divide the preceding estimate by \(m(t)\geq1\) and use the error
bound. The result is exactly
\[
 0\leq1-EN/m(t)
 \leq e^{-u}\{1+(u+M)/d_M\}.
\]
The right side tends to zero independently of \(t\).

For the actual-error version, use the same restart estimate and split
the wrong-exit overshoot \(v=-S_N\) at any \(K\geq1\).
On \(\{v\leq K\}\), its contribution to
\(E[{\bf1}_W(v+M)]\) is at most \(p(K+M)\).
On \(\{v>K\}\), monotonicity of \((v+M)e^{-v}\) gives contribution
at most \((K+M)e^{-K}E[{\bf1}_W e^v]\leq(K+M)e^{-K}\).
This argument does not require \(K\geq u\).
Substitute the resulting bound into
\(m(t)-EN\leq p\,m(t)+E[{\bf1}_W(v+M)]/d_M\), and divide by
\(m(t)\geq1\). This proves the displayed actual-error bound.
For \(0<p\leq e^{-1}\), \(K=\log(1/p)\geq1\) gives its stated form;
\(p\log(1/p)\to0\). If \(p=0\), the restart expectation is zero
and \(EN=m(t)\) exactly.
\end{proof}
\end{annexitem}

\begin{annexitem}{annex:p15.robust.endpoint.closure}{thm:p15-robust-endpoint-closure}{The full neighborhood supremum and every entropy regime}
\begin{proof}
Fix \(i\). The score order of the unchanged neighborhoods and the
proved inverse-CDF construction couple \(Y_i\) under any fixed
\(R\in\mathcal H_i\) above \(Z_i\).
Use one infinite sequence of independent uniforms to carry out the
coupling at every coordinate. The coordinatewise order then holds
simultaneously at all finite times. The negative parts are integrable
by the directional lemma, so no negative infinite increment occurs.
The one-sided hitting time for \(R\) is consequently at most the
one-sided hitting time for \(Q_i\), pathwise in this coupling.
The two-boundary stopping time is at most its own correct one-sided
hitting time. Hence
\[
 E_RN_{a,b}\leq m_i(t_i)
\]
for every member of the entire neighborhood, with the same bound.
It is finite by the truncation theorem.
Since \(Q_i\in\mathcal H_i\), its expectation is a lower bound on
the supremum.

Under \(Q_i\), the directional lemma supplies
\(Ee^{-Z_i}\leq1\) and positive truncated drift.
The two-boundary comparison at \((t_i,u_i)\) gives the remaining
lower bound in the displayed theorem. Once \(\epsilon_i(u_i)<1\),
in particular
\[
 1\leq
 \frac{\sup_{R\in\mathcal H_i}E_RN_{a,b}}{E_{Q_i}N_{a,b}}
 \leq\frac1{1-\epsilon_i(u_i)}.
\]
Let \(u_i\to\infty\). This proves the ratio conclusion, uniformly in
the other positive boundary, without a log-moment assumption.
The actual-error version of the preceding theorem likewise gives
\(E_{Q_i}N/m_i(t_i)\geq1-\delta_{M_i}(p_i)\).
The same upper coupling and membership bounds then squeeze the
worst-case ratio to one whenever \(p_i\to0\), with no requirement
that either boundary diverge. These are precisely the full
vanishing-error quantifiers in the source's least-favorability claim.
The passage limits in \cref{thm:p15-one-sided-truncation} then give
each asserted finite-entropy, infinite-entropy and singular-mass
classification as \(t_i,u_i\to\infty\).

For a fixed \(R\), the directional lemma also gives
\(E_RY_i^-\leq1/e\), \(E_RY_i>0\), and \(E_Re^{-Y_i}\leq1\).
Apply the same one-sided and two-boundary arguments directly under
that law to obtain its individual directional formula.
If \(E_RY_i=D_i<\infty\), both coupled scores are integrable and
their nonnegative difference has expectation zero. It therefore
vanishes almost surely. The score distributions agree, so their
independent sequences and stopping-time laws agree at every boundary.
This reasoning uses finite expectations; it does not subtract
infinity from infinity.
\end{proof}
\end{annexitem}

\begin{annexitem}{annex:p15.stopped.binary.information}{thm:p15-stopped-binary-information}{Stopped-history densities, Wald's identity and binary information}
\begin{proof}
It suffices to prove the null assertion; interchange the laws and
decisions for the other. If \(E_{Q_0}T=\infty\), the lower bound is
immediate. Otherwise assume that expectation finite.
Finite \(D_0\) implies \(Q_0\ll Q_1\) and integrability of
\(Z=\log(dQ_0/dQ_1)\) under \(Q_0\), by the directional lemma.

Let the independent random seed have probability law \(\rho\), with
a one-point space in the absence of randomization.
Consider the disjoint union over \(n\geq1\) of finite stopped-history
spaces \(\mathsf U\times\mathsf X^n\).
Use the sigma-finite reference measure which on component \(n\)
is \(\rho\times\mu^{\otimes n}\), \(\mu=Q_0+Q_1\).
The finite-history stopping rule gives a measurable event \(A_n\)
on that component. Its two stopped-history densities are
\[
 g_i(n,u,x_1,\ldots,x_n)
       ={\bf1}_{A_n}(u,x_1,\ldots,x_n)\prod_{j=1}^nq_i(x_j).
\]
Their integrals over the disjoint union are one because stopping
is almost surely finite under both laws.
On positive \(g_0\), their ratio is the product of the individual
likelihood ratios; the common seed law cancels.
Consequently its log is \(\sum_{j=1}^T Z_j\).

The exact integrability calculation is
\[
 \sum_{j\geq1}E_{Q_0}[|Z_j|{\bf1}_{\{T\geq j\}}]
 =E_{Q_0}|Z|\,E_{Q_0}T<\infty.
\]
The event \(\{T\geq j\}\) is determined by the seed and observations
before \(j\), independent of \(Z_j\). Signed summation is therefore
justified and gives
\[
 D(\text{stopped }Q_0\Vert\text{stopped }Q_1)
 =E_{Q_0}\sum_{j=1}^TZ_j=D_0E_{Q_0}T.
\]
No bounded-overshoot or optional-stopping theorem has been assumed.

For completeness, the decision information inequality follows by
partitioning the stopped space into its two decision events.
If such an event has probabilities \(a,b>0\), conditioning by explicit
restriction and normalization gives
\[
 \int_E g_0\log(g_0/g_1)
 =a\log(a/b)+aD(Q_0(\,\cdot\mid E)\Vert Q_1(\,\cdot\mid E))
 \geq a\log(a/b).
\]
Nonnegativity of the conditional relative entropy is the log inequality
proved in the directional lemma, also valid when the two conditional
laws agree. An event of \(a=0\) contributes zero; if \(a>0=b\),
the information is infinite and the inequality retains that meaning.
Sum over the decision events. Writing
\(x=Q_0(\Delta=1)\), \(y=Q_1(\Delta=1)\), the result is
\[
 D_0E_{Q_0}T\geq
 x\log(x/y)+(1-x)\log((1-x)/(1-y)).
\]
Here \(x\leq\alpha<1-\beta\leq y\).
For \(0<x<y<1\), differentiating this binary expression gives
\[
 \partial_x=\log\frac{x(1-y)}{y(1-x)}<0,\qquad
 \partial_y=\frac{y-x}{y(1-y)}>0.
\]
Its minimum on the permitted rectangle is therefore attained at
\(x=\alpha,y=1-\beta\). Endpoint cases follow by continuous or
infinite limits, with \(0\log0=0\).
Substitution proves the null bound. Interchanging the two labels
gives the alternative expression exactly as stated.
\end{proof}
\end{annexitem}

\begin{annexitem}{annex:p15.sequential.leading.minimax}{thm:p15-sequential-leading-minimax}{Feasible likelihood boundaries and the singular information limit}
\begin{proof}
For every \(R\in\mathcal H_i\), the oriented score has positive
extended mean, integrable negative part and \(E_Re^{-Y_i}\leq1\).
The one-sided proof ensures finite expected two-boundary stopping;
the stopped exponential argument gives wrong-exit probability at most
\(e^{-u_i}\). With \(a=|\log\beta|,b=|\log\alpha|\), these are exactly
the required robust error bounds. Thus this particular test is
admissible for the infima defining \(\mathcal M_i\).

If \(D_i<\infty\), every robustly feasible procedure is feasible at
the least-favorable simple pair. The stopped binary-information
bound is therefore a lower bound on its worst expected cost and
on the infimum of those costs.
For the null, its numerator can be written
\[
 (1-\alpha)|\log\beta|+
 (1-\alpha)\log(1-\alpha)+\alpha\log\alpha
                  -\alpha\log(1-\beta).
\]
The two middle terms are bounded in absolute value, the last tends
to zero, and \(|\log\beta|\to\infty\); its ratio to
\(|\log\beta|\) therefore tends to one.
The alternative expression has the same property relative to
\(|\log\alpha|\).
The upper cost of the chosen likelihood test has the matching limit
by \cref{thm:p15-robust-endpoint-closure}. This squeezes the minimax
cost and proves both finite-entropy assertions.

We next give an endpoint lower bound which remains informative with
infinite entropy. Fix \(i\), let \(j=1-i\), and let
\(C_i=\{q_i>0=q_j\}\), with \(Q_i(C_i)=s_i\).
Let \(T_\infty\) be the first observation in \(C_i\).
For any fixed integer \(k\), decompose the event \(\{T\leq k\}\):
it either makes the wrong decision under \(Q_i\), or makes the correct
decision after seeing \(C_i\), or makes that decision at some
time \(r\leq k\) without having seen \(C_i\).
The first probability is at most the allowed error at \(Q_i\);
the second is at most
\(\Pr_{Q_i}(T_\infty\leq k)=1-(1-s_i)^k\).

For the third, fix \(r\). The \(Q_i\) law of \(r\) observations
restricted to no visit to \(C_i\) is absolutely continuous with
respect to \(Q_j^{\otimes r}\). Its density is
\[
 L_r=
 \begin{cases}
 \displaystyle\prod_{\ell=1}^r(q_i/q_j)(X_\ell),
                  &q_iq_j>0\text{ at every coordinate},\\
 0,&\text{otherwise}.
 \end{cases}
\]
with arbitrary zero assignments on null exceptional sets.
It is nonnegative and integrable, with integral \((1-s_i)^r\leq1\).
The common independent seed does not change this density.
If \(E_r\) is the event of stopping at \(r\) and deciding \(i\),
its probability under \(Q_j\) is at most the other allowed error,
denoted \(\delta_j\). For any finite \(K>0\),
\[
 Q_i(E_r,\text{ no }C_i)
 \leq K\delta_j+
       \int_{\{L_r>K\}}L_r\,dQ_j^{\otimes r}.
\]
First let the allowed errors tend to zero, then \(K\to\infty\).
Integrability makes the tail integral vanish.
For the finitely many \(r\leq k\), this convergence is uniform over
all competing procedures satisfying the error constraints, because
the bound depends only on those constraints and the fixed laws.
It follows that
\[
 \liminf_{\alpha,\beta\to0}\ 
 \inf_{\text{admissible procedures}}\Pr_{Q_i}(T>k)
                      \geq(1-s_i)^k .
\]
Sum this bound over finitely many \(k\geq0\) using
\(ET=\sum_{k\geq0}\Pr(T>k)\).
If \(s_i>0\), increase the finite truncation of the sum to obtain
\(\liminf\mathcal M_i\geq\sum_{k\geq0}(1-s_i)^k=1/s_i\).
The robust likelihood test has worst expected cost tending to that
same constant by the endpoint theorem, proving its minimax limit
and cost ratio.
If \(s_i=0\), the same finite sums can be arbitrarily large, so
\(\mathcal M_i\to\infty\).
When also \(D_i=\infty\), the chosen test has cost \(o(t_i)\);
the nonnegative minimax cost lies below it, so also
\(\mathcal M_i/t_i\to0\).
These bounds do not determine the ratio of two diverging sublinear
costs, which is why no such additional assertion was made.
\end{proof}
\end{annexitem}

\begin{annexitem}{annex:p15.infinite.entropy.finite.scores}{thm:p15-infinite-entropy-finite-scores}{A countable full-support law and its exact truncated drift}
\begin{proof}
The geometric series gives
\(\sum_{j\geq1}2^{-(j+1)}=1/2\), so \(P\) is a probability.
Also \(0<k<(1/2)e^{-2}<1/2\), making \(Q(0)=1-k>1/2\)
and all of its masses positive. Its total mass is one by definition.
The likelihood ratios are \(P(0)/Q(0)=1/\{2(1-k)\}\) at zero
and \(\exp(2^j)\) at \(j\geq1\). This gives the displayed score and
\(0<L<\log2\).
Its negative expectation is \(L/2<\infty\), while the positive
expectation is
\[
 \sum_{j\geq1}2^{-(j+1)}2^j=\sum_{j\geq1}\frac12=\infty.
\]
Thus the relative entropy is infinite and there are no infinite-score
atoms.

For \(M=2^m\), the terms \(j\leq m\) contribute \(m/2\).
The remaining truncated contribution is
\[
 2^m\sum_{j=m+1}^\infty2^{-(j+1)}=\frac12.
\]
Subtract the negative contribution \(L/2\) to get
\(d_M=(m+1-L)/2>0\). At \(t=2^m=M\), the truncation theorem gives
\[
 m_P(t)\leq(t+M)/d_M=\frac{4\,2^m}{m+1-L}.
\]
The absence of infinite scores makes \(m_P(t)\to\infty\), whereas
its infinite entropy gives \(m_P(t)/t\to0\).
For the reverse entropy, the score at zero is \(L\) and the score
at \(j\geq1\) is \(-2^j\). Multiplication by the displayed \(Q\)
masses gives \((1-k)L-\tfrac12\sum_{j\geq1}\exp(-2^j)\).
The series converges because \(2^j\geq j+1\).
Strict positivity follows from the directional entropy lemma and
the distinctness of the two laws.
Zero-radius Huber neighborhoods contain just the respective original
laws and are separated because \(P\ne Q\). The general two-boundary
comparison applies and transfers both conclusions as the two
boundaries diverge.
\end{proof}
\end{annexitem}

\begin{annexitem}{annex:p15.singular.contamination.stopping}{thm:p15-singular-contamination-stopping}{A clipped singular score and two source qualifications}
\begin{proof}
The null class is the singleton \(\delta_0\).
An alternative has the form \((4/5)P_1+(1/5)H\), so its mass \(p\)
at one ranges over \([2/5,3/5]\).
The least-favorable choice is \(H=\delta_0\), giving
\(Q_1(1)=2/5,Q_1(0)=3/5\).
This also follows directly from the source normalization:
with \(\varepsilon_0=0\), \(v=0\), and
\[
 F(u)=\tfrac12+u\quad(u>1/2),\qquad
 F(u)=1/(1-\varepsilon_1)=5/4,
\]
so \(u=3/4\), \(b=4/5\), and the lower score \(bu=3/5\).
The upper score is infinite on the alternative-only point one.

At zero, the log score is \(-c=-\log(5/3)\).
In the null law every observation is zero, so the lower boundary
is reached at exactly \(r=\lceil a/c\rceil\) observations.
There is no null error.
Under an alternative with mass \(p\) at one, the first observation
equal to one causes immediate upper exit, regardless of the finite
upper boundary. If no such observation has occurred by time \(r\),
the test exits at the lower boundary.
Thus \(N=\min(T_p,r)\), where \(T_p\) is geometric with success
probability \(p\), and
\[
 E_RN=\sum_{j=0}^{r-1}(1-p)^j=\frac{1-(1-p)^r}{p}.
\]
Every term is nonincreasing in \(p\), so the worst expectation
over the full alternative class occurs at \(p=2/5\).
Its error is the probability of \(r\) consecutive zeros,
\((3/5)^r\). At \(r=3\), the expectation is
\((1-27/125)/(2/5)=49/25\).

Under \(Q_0\), the directed score is the constant \(c\), giving
\(D_0=c\). Under \(Q_1\), it equals \(+\infty\) with probability
\(2/5\), so \(D_1=\infty\) and \(s_1=2/5\).
The exact expectation tends to \(1/s_1=5/2\).
The allowed choice \(H=P_1\) gives \(R=P_1\), with \(p=1/2\).
Its mean directional score is also infinite, but its expectation
is \(2(1-2^{-r})\), tending to two. This explicitly disproves an
extension of the source's equal-mean argument to two infinite means.

Finally let \(a\to\infty\) with \(b>0\) fixed.
The null error stays zero and the alternative error tends to zero,
although the upper boundary does not diverge. Thus the source's
identification of vanishing errors with both boundaries tending to
infinity requires additional nonsingularity conditions.
The actual-error theorem avoids that invalid implication, and its
infinite-entropy conclusions never divide by infinity inside an
asymptotic equivalence.
\end{proof}
\end{annexitem}
